\documentclass[12pt,oneside]{book}

% packages
\usepackage{amssymb, amsmath, amsthm}

% these make top, right, bottom margins confirming to the requirements

\usepackage[a4paper,bindingoffset=0.0in,%
left=1.5in,right=1in,top=1in,bottom=1in,%
footskip=.25in]{geometry}

\usepackage{setspace}
\usepackage{blindtext}

% small stuff
\usepackage{amsfonts}
\usepackage{amscd}
\usepackage{mathbbol}
\usepackage{hyperref}
\urlstyle{same}

\usepackage{footnote}

\hypersetup{colorlinks,linkcolor={black},citecolor={black},urlcolor={black}}  

\usepackage{microtype}
% \usepackage{showkeys}  
% uncomment this when editing cross-references
\numberwithin{equation}{section}

% this makes the chapters start with a 2" margin
\makeatletter
\renewcommand*\@makechapterhead[1]{%
	\vspace*{0.75in}%
	{\parindent \z@ \raggedright \normalfont
		\ifnum \c@secnumdepth >\m@ne
		\huge\bfseries \@chapapp\space \thechapter
		\par\nobreak
		\vskip 20\p@
		\fi
		\interlinepenalty\@M
		\Huge \bfseries #1\par\nobreak
		\vskip 40\p@
}}
\renewcommand*\@makeschapterhead[1]{%
	\vspace*{0.75in}%
	{\parindent \z@ \raggedright
		\normalfont
		\interlinepenalty\@M
		\Huge \bfseries  #1\par\nobreak
		\vskip 40\p@
}}
\makeatother

% theorem style


\theoremstyle{plain}

\newtheorem{theorem}[subsection]{Theorem}
\newtheorem{proposition}[subsection]{Proposition}
\newtheorem{lemma}[subsection]{Lemma}
\newtheorem{corollary}[subsection]{Corollary}
\newtheorem{conjecture}[subsection]{Conjecture}
\newtheorem{principle}[subsection]{Principle}
\newtheorem{claim}[subsection]{Claim}
%\newtheorem*{main-theorem-1}{Main Theorem}
%\newtheorem*{nil-to-local-repeat}{Proposition \ref{nil-to-local}}
%\newtheorem*{restate-tech-thm}{Theorem \ref{mainthm-tech}}

\theoremstyle{definition}

\newtheorem{definition}[subsection]{Definition}
\newtheorem{problem}{Problem}

\theoremstyle{remark}

\newtheorem{example}{Example}
\newtheorem{examples}{Examples}
\newtheorem*{remark}{Remark}
\newtheorem*{remarks}{Remarks}

%I prefer slanted leq and geq symbols
\renewcommand{\leq}{\leqslant}
\renewcommand{\geq}{\geqslant}

%This is all so I could finish a proof with an equation if need be.
\newsavebox{\proofbox}
\savebox{\proofbox}{\begin{picture}(7,7)%
	\put(0,0){\framebox(7,7){}}\end{picture}}
\def\boxeq{\tag*{\usebox{\proofbox}}}

%These functions are here because I find modular arithmetic difficult in LaTeX.
\newcommand{\md}[1]{\ensuremath{(\operatorname{mod}\, #1)}} % \operatorname automagically uses the right font and size

%\def\proof{\noindent\textit{Proof. }}
\def\endproof{\hfill{\usebox{\proofbox}}}
\def\E{\mathbb{E}}
\def\so{\mathfrak{so}}
\def\sp{\mathfrak{sp}}
\def\Z{\mathbb{Z}}
\def\R{\mathbb{R}}
\def\T{\mathbb{T}}
\def\C{\mathbb{C}}
\def\N{\mathbb{N}}
\def\P{\mathbb{P}}
\def\F{\mathbb{F}}
\def\Q{\mathbb{Q}}
\def\b{{\mathbf b}}
\def\A{{\mathbf A}}
\def\g{{\mathfrak g}}
\def\hp{{\mathfrak{hp}}}
\def\D{\mathcal{D}}
\def\re{\operatorname{Re}}
\def\O{\mathcal{O}}

\newcommand\DD{{\frac{\partial}{\partial z}}}
\newcommand\DDbar{{\frac{\partial}{\partial \overline{z}}}}

\newcommand\SPEC{{\textrm{Spec}}}
\newcommand\MM{{\mathscr M}}
\newcommand\OO{{\mathscr O}}
\newcommand\LL{{\mathcal L}}
\newcommand\PP{{\mathbf P}}
\newcommand\eps{{\varepsilon}}
\newcommand\supp{\mathop{\rm supp}}
\newcommand\Riem{{\operatorname{Riem}}}
\newcommand\dist{{\operatorname{dist}}}
\newcommand\Hess{{\operatorname{Hess}}}
\newcommand\Ric{{\operatorname{Ric}}}
\newcommand\Vol{{\operatorname{Vol}}}
\renewcommand\min{{\operatorname{min}}}

\parindent 0mm
\parskip   5mm % this should make \ni and  obsolete, hopefully
%\newenvironment{proof}{\noindent {\bf Proof} }{\endprf\par}
\def \endprf{\hfill  {\vrule height6pt width6pt depth0pt}\medskip}

% add a customized "itemize" operator. 

\def\+{+}

\makeatletter
\catcode`\ =12\let\@nl@space= \catcode`\ =10
\newcount\@nl@rlevel
\newcount\@nl@llevel
\@nl@llevel=-1

\def\@nl{%
	\catcode`\ =12
	\global\@nl@rlevel=0
	\futurelet\@nl@store\@nl@%
}
\def\@nl@gobble#1{\futurelet\@nl@store\@nl@}
\def\@nl@enditemize{
	\ifnum\the\@nl@rlevel<\the\@nl@llevel%
\end{itemize}%
\egroup%
\expandafter\@nl@enditemize%
\else%
\ifnum\the\@nl@rlevel=\the\@nl@llevel\else%
\errmessage{Error: inconsistent identation}
\fi%
\fi%
}
\def\@nl@{%
\ifx\@nl@store\@nl@space%
\global\advance\@nl@rlevel by 1
\expandafter\@nl@gobble%
\else%
\catcode`\ =10
\ifx\@nl@store+%
\ifnum\the\@nl@rlevel>\the\@nl@llevel%
\bgroup%
\@nl@llevel=\the\@nl@rlevel
\begin{itemize}%
	\fi%
	\@nl@enditemize%
	\item \expandafter\expandafter\expandafter\@gobble%
	\else%
	\ifx\@nl@store\@nl%
	\global\@nl@rlevel=-1\relax\@nl@enditemize\par
	\else\space\fi%
	\fi%
	\fi%
}

\catcode`\^^M=\active%
\AtBeginDocument{%
	\catcode`\^^M=\active%
	\let^^M=\@nl%
}%
\catcode`\^^M=5
\makeatother
%------------------------------------------------




% The next lines are personal shortcuts and settings.  Feel free
% to delete those you don't want.

%  theorems, corollaries, etc.


% doesn't count introduction when numbering theorems, etc.
\setcounter{chapter}{0}

% makes page numbers appear
\pagestyle{plain} 


\begin{document}

% makes the page numbers roman numerals, doesn't count
% these pages in the table of contents
\frontmatter

\thispagestyle{empty}

\vbox to 1truein{}

\centerline{THESIS OR DISSERTATION TITLE IS CENTERED AND IN ALL CAPS}



\vskip 200pt

\centerline{BY}
\vskip 10pt

\centerline{XXX}
\vskip 10pt

\centerline{BA, YYY College, 2222}
\centerline{MA, ZZZ University, 2224}

\vskip 215pt

\centerline{DISSERTATION}
\vskip 10pt

\centerline{Submitted in partial fulfillment of the requirements for}
\centerline{the degree of Doctor of Philosophy in EEE}
\centerline{in the Graduate School of}
\centerline{Binghamton University}
\centerline{State University of New York}
\centerline{2229}

\newpage

\thispagestyle{empty}

\vbox to 9.0truein{}

\centerline{\copyright\ Copyright by XXX 2229}

\

\centerline{All Rights Reserved}

\newpage

{\baselineskip = 10pt

\vbox to 2.0truein{}


\vskip 295pt

\centerline{Accepted in partial fulfillment of the requirements for}
\centerline{the degree of Name of Degree in Major}
\centerline{in the Graduate School of}
\centerline{Binghamton University}
\centerline{State University of New York}
\centerline{2229}

\

\centerline{April 1st, 2229}

\

\centerline{AAA, Chair}
\centerline{Department of EEE, Binghamton University}


\

\centerline{BBB, Faculty Advisor}
\centerline{Department of EEE, Binghamton University}

\

\centerline{CCC , Member}
\centerline{Department of EEE, Binghamton University}

\

\centerline{DDD, Outside Examiner}
\centerline{Department of FFF, GGG University}
}

\chapter*{Abstract}
\begin{doublespace}
Blah Blah Blah
\end{doublespace}


\newpage

\chapter*{Dedication}
\bigskip

\centerline{Dedicated to HHH}

\newpage

\chapter*{Acknowledgements}
Blah Blah Blah


\newpage

\tableofcontents

% Changes page numbers to regular numbers, resets the counter
\mainmatter

% This gives 11pt font with 20pt spacing, text from here should be double spaced
\fontsize{11}{20pt} \selectfont  

% \include puts in the .tex file with the given name
% make sure that these files don't have any preamble material

\begin{doublespacing}
	
\include{introduction}

\include{chapter1}

\include{chapter2}

\include{chapter3}

\include{conclusions}

\end{doublespacing}

% add a new chapter without a chapter # for the references
\addcontentsline{toc}{chapter}{Bibliography} 

\begin{thebibliography}{9}

\bibitem[Voisin 2002]{Voison}Claire Voisin, \emph{Hodge theory and complex algebraic geometry I, II},  Cambridge
University Press, New York, 2002, ix+322 pp., £65, ISBN 0-521-80260-1 (Vol. I);
2003, ix+351 pp., £65, ISBN 0-521-80283-0 (Vol. II)

\end{thebibliography}



\end{document}
